% answers_short.tex —— 2014 年·北京卷·文科 · 简短答案速查
% 由 tools/make_short_answers.py 自 answers.tex 的 \gksolhead/\gkanswer 提取，
% 并与 2014/2014_数学真题_北京文_答案要点.md 逐题比对后生成，勿手改（改源 MD 后重跑）。

\begin{gknotice}
\textbf{使用说明}：本节只给最终答案便于快速核对，完整推导见解析版。答案为编写者独立推导，\textbf{非官方标准答案}。
\end{gknotice}

\gkpart{选择题}
\begin{center}\begin{tabular}{ccccccccc}
\toprule
题号 & 1 & 2 & 3 & 4 & 5 & 6 & 7 & 8 \\
\midrule
答案 & C & B & A & C & D & C & B & B \\
\bottomrule
\end{tabular}\end{center}

\gkpart{填空题}
\begin{itemize}[leftmargin=2.6em]
  \item[\textbf{9.}] $2$。$(x+\mathrm i)\mathrm i=-1+x\mathrm i=-1+2\mathrm i$，比较虚部得 $x=2$。
  \item[\textbf{10.}] $x^2-y^2=1$。$c=\sqrt2$，$a=1$，故 $b^2=c^2-a^2=1$。
  \item[\textbf{11.}] $2\sqrt2$。还原得 $AB=2$、$BC$ 方向深 $1$、高 $2$，最长棱 $\sqrt{2^2+2^2}=2\sqrt2$。
  \item[\textbf{12.}] $c=2$；$\sin A=\dfrac{\sqrt{15}}8$。$c^2=1+4-2\cdot1\cdot2\cdot\dfrac14=4$，$\cos A=\dfrac78$。
  \item[\textbf{13.}] $1$。可行域顶点为 $(0,1)$，$(2,1)$，$(1,0)$，$z$ 值分别为 $1$，$2\sqrt3+1$，$\sqrt3$。
  \item[\textbf{14.}] $42$。工艺师精加工共 $36$ 个工作日，最早第 $6$ 天开工，且该下界可由「先粗加工 $B$ 再粗加工 $A$」达到。
\end{itemize}

\gkpart{解答题}
\begin{itemize}[leftmargin=2.6em]
  \item[\textbf{15.}] (1) $a_n=3n$，$b_n=3n+2^{n-1}$。公差 $d=3$，等比数列 $\{b_n-a_n\}$ 首项 $1$、公比 $2$；(2) $S_n=\dfrac{3n(n+1)}2+2^n-1$。分组求和：等差数列部分与等比数列部分分别求和
  \item[\textbf{16.}] (1) $T=\pi$，$x_0=\dfrac{7\pi}6$，$y_0=3$。最大值点 $x=\dfrac\pi6+k\pi$，图中 $x_0$ 为第二个最大值点；(2) 最大值为 $0$，最小值为 $-3$。$2x+\dfrac\pi6\in\left[-\dfrac{5\pi}6,0\right]$，$\sin$ 在 $u=-\dfrac\pi2$ 处取 $-1$，端点处最大为 $\sin0=0$
  \item[\textbf{17.}] (1) 见解析版；(2) 见解析版；(3) $\dfrac{\sqrt3}3$。$AB=\sqrt3$，$S_{\triangle ABC}=\dfrac{\sqrt3}2$，高为 $AA_1=2$，$V=\dfrac13\cdot\dfrac{\sqrt3}2\cdot2$
  \item[\textbf{18.}] (1) $0.9$。少于 $12$ 小时的频数为 $90$，$\dfrac{90}{100}=0.9$；(2) $a=0.085$，$b=0.125$。矩形高为 $\dfrac{\text{\text{频数}}}{200}$，第 $3$ 组 $17$、第 $5$ 组 $25$；(3) 第 $4$ 组。中点值加权平均得 $\bar x=7.68\in[6,8)$
  \item[\textbf{19.}] (1) $\dfrac{\sqrt2}2$。$a=2$，$b=\sqrt2$，$c=\sqrt2$；(2) $2\sqrt2$。$|AB|^2=|OA|^2+|OB|^2=\dfrac{(u+2)^2}{u}\geqslant 8$，其中 $u=2-y^2\in(0,2]$，$u=2$ 时取等号
  \item[\textbf{20.}] (1) $\sqrt2$。$f'(x)=3(2x^2-1)$，极大值 $f\left(-\dfrac{\sqrt2}2\right)=\sqrt2$ 大于端点值 $f(-2)=-10$、$f(1)=-1$；(2) $(-3,-1)$。切点横坐标满足 $-4x^3+6x^2-3=t$，该函数极小值 $-3$、极大值 $-1$，三根当且仅当 $-3<t<-1$；(3) 过 $A$ 有 $3$ 条，过 $B$ 有 $2$ 条，过 $C$ 有 $1$ 条。对应方程 $4x^3+6x^2-1=0$（三个实根）、$(x+1)(x-2)^2=0$、$4x^3+2=0$（一个实根）
\end{itemize}
